% Review 2010-08-13: NASA-Bild kontrovers
% \begin{savequote}
% \includegraphics[width=5cm]{./images/nasa-pd/157684main_image_feature_656_ys_full-00800.jpg}
% 
% Bild: NASA
% \end{savequote}
\chapter
[\CO{[GYN]} \emph{Spezielle Patientengruppe:} Frauen und Schwangere]
{\CC{[GYN]} \emph{Spezielle Patientengruppe:} Frauen und Schwangere}

\ChapterIntro
{%
~~~
}
{%
\begin{description}
  \item[Maintainer:] Sebastian Gabriel
  \item[Autoren:] Diverse
  \item[Reviewer:] Standard-Reviewprozess
  \item[Version:] {\DocVersion} ({\DocVersionInfo})
  \item[SHA1:] {\DocGitCommit}
\end{description}

{\TxTeilkapitelAASS}
}


\clearpage % TEMP temp clearpage
\section{Schwangerschaft}
\HeadingSub{\Lang{Term.} {Gestation}}

\subsection{Allgemeines}

\Layout{0}{Exkurs: Weiblicher Zyklus}{%
Der Begriff Zyklus oder Periode beschreibt in diesem
Zusammenhang den Vorgang des rund alle \EE{28 Tage}
stattfindenden Reifens mindestens einer Eizelle, dem
Eisprung und den damit im Zusammenhang stehenden
Veränderungen der Gebärmutterschleimhaut. Der Zyklus ist
unter
\FullPrettyRef{topic:weiblicher-zyklus}, beschrieben.
}
{\FullPrettyRef{topic:weiblicher-zyklus}}
{}{}{}









\Layout{0}{Schwangerschaftsdauer}{%
Die normale Schwangerschaft dauert \EE{40\,$\pm$\,2
Wochen}. Das sind rund \EE{9 Kalendermonate} \Z{oder 10
Lunarmonate\footnote{\Z{\Lang{lat.} luna: Mond.}}}. Der
Fortschritt der Schwangerschaft wird in
\T{Schwangerschaftswochen}, Abk. \T{SSW}, angegeben. Häufig
wird auch die Unterteilung in drei
\E{Trimester}\ZFootnote{\Lang{lat.} tri: drei; \Lang{lat.}
mensis: Monat; Trimenon, Trimester: dreimonatig} zu je 3
Kalendermonaten verwendet.
Von
einer \T{Frühgeburt} spricht man wenn die Geburt
\E{bis zur vollendeten 37. SSW}
erfolgt. Von einer \E{Übertragung} spricht man ab der
\E{42.~SSW}.
}{%
\begin{itemize}
  \item 40 $\pm$ 2 Schwangerschaftswochen (SSW)
  \item 3 Trimester
  \item Frühgeburt: bis 37. SSW
  \item Übertragung: > 42. SSW
\end{itemize}
}{}{}{}


\Layout{0}{Schwangerschaftsverlauf}
{%
Nach dem Eisprung ist die Eizelle \VonBis{8}{10} Stunden
befruchtungsfähig. Sie benötigt für die Wanderung Eierstock
\Sign{→} Eileiter \Sign{→} Uterus (Gebärmutter)
\VonBis{4}{6} Tage. Wurde die Eizelle befruchtet, kommt es
zur \EE{Einnistung in der Gebärmutterschleimhaut}. Dort
bildet sich der \T{Mutterkuchen} (\T{Plazenta}), welcher die
Frucht\footnote{Mit dem Begriff \Speech{\enquote*{Frucht}}
ist das Kind gemeint. Die  Frucht hat je nach
Schwangerschaftsfortschritt verschiedene Namen: Man spricht
vom \I{Embryo} bis zur 12. SSW, danach vom \I{Fetus}.} mit
Nährstoffen und Sauerstoff versorgt.
Weiters bildet sich um die Frucht eine \T{Fruchtblase},
welche mit \T{Fruchtwasser} gefüllt ist und sie so ideal vor
äußeren Erschütterungen schützt.

Zwischen der 18. und der 20. SSW kann die Mutter die ersten
Kindesbewegungen spüren. \EE{Ab der 24.
SSW gilt das Kind als überlebensfähig} und ggfs. bei einer
Frühgeburt als reanimationspflichtig.
}
{%
\begin{itemize}
  \item Eierstock \Sign{→} Eileiter \Sign{→} Gebärmutter
  \item Befruchtung
  \item Einnistung in Gebärmutterschleimhaut
  \item Mutterkuchen (Plazenta)
  \item Fruchtblase, Fruchtwasser
  \item Ab 24. SSW überlebensfähig
\end{itemize}

}
{}
{}
{}



\begin{PictureSeriesWide}{}{Bilderserie: Schwangerschaft}
\PictureSeriesRowWide{3}
{./images/wikipedia-pd/Nabelschnur.png}
{Die Lage des Fötus in der Gebärmutter.
\SourcePicture{Gray's Anatomy, 20th ed., PD, Copyright
expired.}}
{empty}
{}
{empty}
{}
{empty}
{}
\end{PictureSeriesWide}




\subsection{Notfälle -- Frühschwangerschaft}
\label{topic:notfall-fruehschwangerschaft}

\subsubsection{Fehlgeburt}
\HeadingSub{\Lang{Term.} \T{Abortus}}
%
\index{Fehlgeburt}%
\index{Abort}%
\label{topic:fehlgeburt}%
\label{topic:abort}%

\Definition{Beendigung einer Schwangerschaft bei einem
Kindesgewicht von unter 500\Gramm* und fehlenden Lebenszeichen.
\cite{HaagGyn2007,ForMed2}}


\Layout{0}{\TxBeschreibung}{%
Aufgrund verschiedener Ursachen kann es zum Absterben der
Frucht mit einer nachfolgenden Fehlgeburt (Ausstoßung)
kommen. Wenn dies vor dem Erreichen von 500\Gramm*
geschieht, spricht man von einer \T{Fehlgeburt}.
}{
\begin{itemize}
  \item Frühe Fehlgeburt\,/\,Fruchtabgang
  \item Oft unbemerkt
\end{itemize}
}{}{}{}


\Layout{0}{\TxSymptome}{%
Die Symptome sind vor allem davon abhängig, wie weit die
Schwangerschaft fortgeschritten ist bzw. die Frucht
entwickelt ist. Ein sehr früher Fruchtabgang wird oft nicht
einmal bemerkt bzw. mit Regelunregelmäßigkeiten verwechselt,
oft weiß die Frau nicht einmal, dass sie schwanger gewesen
ist.  Ein Abort macht sich dann vor allem durch eine
\EE{vaginale Blutung} bemerkbar, die mitunter als
\Speech{\enquote{verstärkte Regelblutung}} fehlgedeutet wird.
Weiters kann es zum \EE{Abgang von Gewebeteilen} und
\EE{Unterbauchschmerzen} kommen.
Eine Fehlgeburt kann jedoch auch zu einem späteren Zeitpunkt
während der Schwangerschaft vor Erreichen der
500\Gramm*-Grenze des Kindes vorkommen.\ZFootnote{Nach dem
Erreichen der 500\Gramm*-Marke spricht man von einer
\E{Totgeburt}; für diese gelten z.\,T. andere rechtliche
Rahmenbedingungen (Kind kann einen Vornamen erhalten,
Beurkundung im Sterbebuch).
\cite{ForMed2}.}
}{
\begin{itemize}
	\item Vaginale Blutung
	\item Abgang von Gewebsteilen
	\item Unterbauchschmerzen
\end{itemize}
}{}{}{}

\Layout{0}{Anmerkung}{% 
Sollten Zweifel bestehen, ob es sich um einen Abort handelt,
sollte die Verdachtsdiagnose \Speech{\enquote{V.\,a. Abort}}
\EE{vermieden} werden, um die Patientin nicht --
möglicherweise unnötig -- zu beunruhigen. Stattdessen sollte
man eine korrekte Symptombeschreibeung als Verdachtsdiagnose
verwenden, z.\,B. \Speech{\enquote{Vaginale Blutung in 5. SSW}}.

Sollten Zweifel hinsichtlich der
Lebensfähigkeit des Kindes bestehen, muss mit einer
Reanimation begonnen werden.

\Cave{Im Zweifel: Reanimation!}
}{% 
Im Zweifel reanimieren! 

Diagnose \enquote{Abort} vermeiden \Sign{→} Symptombeschreibung}{}{}{}

\MassnahmenMK{Spezielle Maßnahmen}{Abort}{m:abort}{
}{}
% \MassnahmenNG{Spezielle Maßnahmen}{Abort}{\label{m:abort}}{
% \begin{itemize}
%   \item Lagerung: sterile Vorlage und Lagerung nach
%   Fritsch
%   % TODO Querverweis zu Lagerungsarten einfügen
%   \item \TxBeiVit: \TxMassVitBes
%   \begin{itemize}
%     \item Lagerung: Situationsgerecht
%   \end{itemize}
%   \item Abgegangene Gewebeteile sammeln
%   \item Transportentscheidung: Abt. f. Gynäkologie und
%   Geburtshilfe
% \end{itemize}
% } 




\subsubsection{Eileiterschwangerschaft}
\label{topic:extrauterine-graviditaet}

\HeadingSub{\Lang{Term.} {Tubaria}, \Lang{Syn.} extrauterine Schwangerschaft}



\Layout{0}{\TxBeschreibung}{%
\index{Eileiterschwangerschaft}%
Wenn sich die befruchtete Eizelle nicht in der Gebärmutter,
sondern schon im Eileiter einnistet, kommt es zu einer
\EE{Eileiterschwangerschaft}
die lebensbedrohlich ist.
Mit zunehmender Größe der Frucht wird der Eileiter gedehnt,
dabei kommt es zu Schmerzen. Schließlich reißt der Eileiter
und die Frau erleidet massive innere Blutungen. Es zeigen
sich dann Anzeichen eines akuten Abdomens und eines
hypovolämischen Schocks.
}{
\begin{itemize}
    \item Einnistung der Frucht in Eileiter (statt Gebärmutter)
    \item Schmerz durch Dehnung des Eileiters
    \item \Sign{→} Zerreißung des Eileiters
    \begin{itemize}
        \item Blutung
        \item Akutes Abdomen
        \item Schock
    \end{itemize}
\end{itemize}
}{}{}{}



\Fazit{Bei Frauen mit akutem Abdomen muss immer eine mögliche
Schwangerschaft abgeklärt werden!}


\MassnahmenMK{Spezielle Maßnahmen}{Eileiterschwangerschaft,
Verdacht}{m:eileiterschwangerschaft-verdacht}{
}{}

\SlideCompensation{2}

\MassnahmenMK{Spezielle Maßnahmen}{Eileiterschwangerschaft,
rupturiert}{m:eileiterschwangerschaft-rupturiert}{%
}{}
% \MassnahmenNG{Spezielle Maßnahmen}{Eileiterschwangerschaft,
% Verdacht}{\label{m:eileiterschwangerschaft-verdacht}}{
% \begin{itemize}
%     \item Transportentscheidung: Abt. f. Gynäkologie
% \end{itemize}
% }
% 
% \MassnahmenNG{Spezielle Maßnahmen}{Eileiterschwangerschaft,
% rupturiert}{\label{eileiterschwangerschaft-rupturiert}}{%
% \begin{itemize}
%   \item \TxMassVitBes
%   \begin{itemize}
%     \item Lagerung: Bauchdeckenentspannend.
%     \item Schockbekämpfung
%     \item \ce{O2}: \VonBis{10}{15}\LPerMin*
%   \end{itemize}
% \end{itemize}
% }


\clearpage % TEMP temp clearpage
\subsection{Notfälle -- Spätschwangerschaft}
\label{topic:notfall-spaetschwangerschaft}

\subsubsection{Vorzeitige Plazentalösung}


\Definition{Ablösung des Mutterkuchens (Plazenta) vor Beendigung der Geburt}


\Layout{0}{\TxBeschreibung}{%
Aufgrund einer Einblutung zwischen Plazenta und Gebärmutter
kann es zu einer vorzeitigen Ablösung der Plazenta mit dann
ungehinderter \EE{Blutung} kommen. Es besteht Schock- bzw.
\EE{Lebensgefahr} für Mutter und Kind durch Blutung der Mutter und
Mangelversorgung des Kindes.
Die Diagnose ist präklinisch nicht sicher zu
stellen. Im Vordergrund steht daher die Behandlung der
Symptome und ein zügiger Transport.
}
{
\begin{itemize}
  \item Blutung
  \item Schock, Lebensgefahr
\end{itemize}
}
{}{}{}


\Layout{0}{\TxSymptome}{%
Die Symptome sind unspezifisch: Im Vordergrund steht der
\EE{Schock} infolge der Blutung. Zusätzlich zu den
Schocksymptomen kann die Frau Symptome eines \EE{akuten
Abdomens} zeigen.

\bigskip

\bigskip
}{%
\begin{itemize}
  \item Unspezifisch:
  \begin{itemize}
    \item Schock
    \item Akutes Abdomen
  \end{itemize}
\end{itemize}
}{}{}{}


\MassnahmenMK{Spezielle Maßnahmen}{Vorzeitige
Plazentalösung}{m:plazentaloesung}{
}{}
% \MassnahmenNG{Spezielle Maßnahmen}{Vorzeitige
% Plazentalösung}{\label{m:plazentaloesung}}{
% \begin{itemize}
%   \item \TxMassVitBes
%   \begin{itemize}
%     \item Lagerung: sterile Vorlage und Lagerung nach Fritsch.
%     \item Schockbekämpfung
%   \end{itemize}
%   \item Keine vaginale Untersuchung oder  Scheidentamponade
%   \item Zügiger Transport mit Voranmeldung
% \end{itemize}
% }


\subsubsection{Vena-cava-Syndrom}
\label{topic:vena-cava-syndrom}

\Layout{2}{\TxBeschreibung}{%
Bei hochschwangeren Frauen kann es vorkommen, dass das
\EE{Kind die untere Hohlvene der Mutter abdrückt}. Dadurch
wird der Blutstrom zum Herzen vermindert und es kann zu
einem \EE{plötzlichen Kollaps der Mutter} und einer
Unterversorgung des Kindes kommen. Durch eine
Linksseitenlagerung wird der Druck, den das Kind ausübt,
vermindert, da die Hohlvene eher rechts der Körpermitte
verläuft.

\Fazit{\EE{Jede hochschwangere Patientin }sollte daher in
Linksseitenlage und nicht in Rückenlage transportiert werden. }
}{

}
{
\begin{itemize}
  \item Kind drückt untere Hohlvene der Mutter ab
  \item Kollaps
\end{itemize}
}
{}
{}

\MassnahmenMK{Spezielle
Maßnahmen}{Vena-cava-Syndrom}{m:vena-cava-syndrom}{
}{}
% \MassnahmenNG{Spezielle
% Maßnahmen}{Vena-cava-Syndrom}{\label{m:vena-cava-syndrom}}{
% 
% \begin{itemize}
%     \item Linksseitenlage, grundsätzlich immer bei
%     hochschwangeren Patientinnen.
% \end{itemize}
% }




\subsubsection{Präeklampsie, EPH-Gestose, HELLP-Syndrom, Eklampsie}
\label{topic:eph-gestose}\label{topic:praeeklampsie}\label{topic:eklampsie}\index{EPH-Gestose}\index{Präeklampsie}\index{Eklampsie}\index{Krampfanfall!eklamptischer}

% FEATURE INHALT: HELLP,Syndrom, Leberspannungsschmerz

\Layout{0}{\TxBeschreibung}{%
Die \T{EPH-Gestose}, auch \T{Präeklampsie} genannt,
ist eine gefürchtete Komplikation der
Spätschwangerschaft, bei der es zu einer
Durchblutungsstörung der Gebärmutter, und in weiterer Folge
zu diversen Komplikationen bei der Mutter kommt. Sie wird
in der Regel im Rahmen der Mutter-Kind-Pass-Untersuchungen
erkannt. Die
Maximalvariante der Präeklampsie ist das \T{HELLP-Syndrom}.
Unbehandelt kann die Präeklampsie und das HELLP-Syndrom zu
einem \E{eklamptischen Krampfanfall} (\T{Eklampsie}, ähnlich
einem \enquote*{normalen} zerebralen Krampfanfall) und in
weiterer Folge zum Tod von Mutter und Kind führen.

\Fazit{Die \emph{Prä}eklampsie (EPH-Gestose) kann zu einem
lebensgefährlichen Krampfanfall (Eklampsie) führen.}
}
{
\begin{itemize}
  \item Schwangerschaftserkrankung
  \item Kann zur Eklampsie (Krampfanfall) führen
\end{itemize}
}
{}
{}
{}



\Layout{0}{\TxSymptome}{
Die Präeklampsie ist charakterisiert durch:
\begin{enumerate}
    \item \EE{E}dema (Ödeme)
    \item \EE{P}roteinuria (Proteinurie = Eiweiß im Harn, Harn schäumt)
    \item \EE{H}ypertonie (RR  ${\uparrow}$) 
    
    Ein Blutdruck
    {\textgreater}\,\RR{140}{90} bei einer Schwangeren ist zu  hoch!
\end{enumerate}
Beim HELLP-Syndrom kommt es zu einem \EE{rechten
Oberbauchschmerz} (Leberdehnungsschmerz).
Bei einer drohenden Eklampsie klagt die Patientin über 
\EE{Kopfschmerzen}, Sehstörungen sowie Übelkeit und
Erbrechen. \cite{Haag2006}

}{
\begin{itemize}
    \item \EE{E}dema (Ödeme)
    \item \EE{P}roteinuria (Proteinurie = Eiweiß im Harn)
    \item \EE{H}ypertonie (RR  ${\uparrow}$)
    \item HELLP: re. Oberbauchschmerz
    \item Drohende Eklampsie: \begin{compactitem}
      \item Kopfschmerzen
      \item Sehstörungen
      \item Übelkeit, Erbrechen
    \end{compactitem}
\end{itemize}
}{}{}{}


\SlideCompensationNoParskip{1}





\MassnahmenMK{Spezielle Maßnahmen}{Verdacht auf
Präeklampsie}{m:praeeklampsie-verdacht}
{
}{}

\SlideCompensation{1}

\MassnahmenMK{Spezielle Maßnahmen}{Eklampsie --
Eklamptischer Krampfanfall}{m:eklampsie}{
Grundsätzlich ist so wie bei jedem anderen Krampfanfall zu
verfahren , vgl. Kap. \ref{topic:krampfanfall}. Bedenke:
Der eklamptische Krampfanfall ist für Mutter und Kind
lebensgefährlich!
}{}
% \MassnahmenNG{Spezielle Maßnahmen}{Verdacht auf
% Präeklampsie}{\label{m:praeeklampsie-verdacht}}
% {
% \begin{itemize}
%     \item schonender Transport
%     \item Überwachung, evtl. \ce{O2}
%     \item Reizabschirmung
%     %\item Notarzt: Krampfdurchbrechung, RR-Senkung
%     %\item (im Vollbild: vorzeitige Beendigung der  Schwangerschaft)
% \end{itemize}
% }
% 
% \MassnahmenNG{Spezielle Maßnahmen}{Eklampsie --
% Eklamptischer Krampfanfall}{\label{m:eklampsie}}{
% Grundsätzlich ist so wie bei jedem anderen Krampfanfall zu verfahren , vgl. Kap. \ref{topic:krampfanfall}. Bedenke:
% Der eklamptische Krampfanfall ist für Mutter und Kind lebensgefährlich!
% }


\subsubsection{Vorzeitiger Fruchtwasserabgang}

\Layout{2}{\TxBeschreibung}{%
Der Fötus schwimmt während  Schwangerschaft in steriler
Flüssigkeit (Fruchtwasser). Bei vorzeitiger Eröffnung der
Fruchtblase ergeben sich für ihn die Gefahr
eines \EE{Nabelschnurvorfalls} oder einer \EE{Infektion}
(Keimbesiedlung der Fruchthöhle). Die Patientin ist an
einer geburtshilflichen Abteilung vorzustellen.
}{}{}{}{}


\MassnahmenMK{Spezielle Maßnahmen}{Vorzeitiger
Fruchtwasserabgang}{m:fruchtwasserabgang}{
}{}



\clearpage
\section{Geburtshilfe}

\subsection{Die Geburt}
\label{topic:geburt}

Als \EE{Beginn der Geburt} zählt:
\begin{itemize}
  \item wenn regelmäßige Wehen in \VonBis{10}{15} Minuten
  Abstand auftreten und dieser Zustand länger als eine
  Stunde dauert \item wenn die Fruchtblase springt und
  Fruchtwasser abgeht
\end{itemize}

Die Geburt verläuft in 3 Phasen: 
\begin{itemize}
  \item \T{Eröffnungsphase}
  \item \T{Austreibungsphase} 
  \item \T{Nachgeburtsphase}
\end{itemize} 


\Layout{0}{Eröffnungsphase}
{%
Die Eröffnungsphase beginnt mit den ersten regelmäßigen
Wehen und endet mit der vollständigen Eröffnung des
Muttermundes auf 10\CM*.
\cite{StriebelAnaesthesie2,BuehlingGynaekologie1,HaagGyn2007}
Durch Zusammenziehen der Gebärmuttermuskulatur
(Wehen) wird das Kind langsam in den Geburtskanal nach unten
gepresst. Dadurch verkürzt sich der Gebärmutterhals
(Zervix), und der Muttermund wird eröffnet. Der
Wehenabstand beträgt ca. \VonBis{2}{10}\Min*.
Am Ende der Eröffnungsphase
kommt es zum Platzen der Fruchtblase (\T{Blasensprung}),
dabei geht Fruchtwasser und etwas blutiges Sekret ab.

\Z{Bei Erstgebärenden dauert die Eröffnungsperiode ca.
\VonBis{6}{9}\Hour*, bei Mehrgebärenden ca.
\VonBis{3}{7}\Hour* \cite{BuehlingGynaekologie1}. Die
tatsächliche Dauer kann sehr unterschiedlich sein.}

\Cave{%
Ab dem Blasensprung darf die Gebärende nicht mehr aufstehen,
sondern muss liegend transportiert werden. Es besteht die
Gefahr eines Nabelschnurvorfalls in den Geburtskanal!}
}{%
\begin{itemize}
  \item Von Beginn regelmäßger Wehen bis vollst. Eröffnung
  des Muttermundes
  \item Wehen im Abstand von \VonBis{2}{10} Minuten
  \item Gebärmutterhals verkürzt sich, Muttermund öffnet sich
  \item Blasensprung mit Fruchtwasserabgang
  \item Ab Blasensprung: \EE{Gefahr eines
  Nabelschnurvorfalles}
  \begin{itemize}
    \item \Ca{Liegender Transport!}
  \end{itemize}
\end{itemize}
}{}{}{}

\Layout{0}{Austreibungsphase}
{%
Die Austreibungsphase beginnt mit der
vollständigen \E{Eröffnung} des Muttermundes \Z{(10\CM*)}
und endet mit der \E{Geburt} des Kindes.
\cite{BuehlingGynaekologie1,HaagGyn2007}
Der Wehenabstand beträgt nun \EE{\VonBis{2}{3}}
Minuten. Die Wehen werden im Verlauf kraftvoller und die
Frau verspürt mit Fortschreiten der Geburt einen
\EE{Pressdrang}(\EE{Presswehen}). Der Kopf des Kindes wird
schließlich nun zwischen den Schamlippen  sichtbar.
Wenn der Kopf geboren ist, folgt der restliche Körper in der
Regel schnell nach. Vorsicht: Das Kind ist durch die
Fruchtschmiere rutschig!

\subparagraph{Der Damm}
\label{topic:dammschutz}
Als \T{Damm} wird das Stück Haut zwischen Scheide und Anus
bezeichnet. Mittels eines
Griffes namens \T{Dammschutz} kann
theoretisch die Wahrscheinlichkeit eines Dammrisses
minimiert werden. Dabei wird die Hand unterhalb der
Scheide angelegt und man versucht das Gewebe zu
stabilisieren. In der Praxis ist dieser Schutz, insbesonders
für darin nicht routiniertes Personal, schwer zu erreichen
\cite{AsboeLehrmeinung2011-20}.

Dennoch ist der Dammschutz nicht unwichtig: Durch das
Pressen wird nicht nur das Kind geboren, sondern in der
Regel auch  \EE{Stuhl} abgesetzt. Durch den Dammschutz, in
Verbindung mit einer Unterlage, kann das Kind \E{vor dem
Stuhl der Mutter geschützt} werden (\T{Stuhlschutz}).

}{%
\begin{itemize}
  \item Von vollständiger Eröffnung des Muttermundes bis
  Geburt des Kindes
  \item Wehen im Abstand von \VonBis{2}{3}\Min*
  \item Presswehen
  \item Dammschutz = Stuhlschutz
  \item Geburt des Kindes
\end{itemize}
}
{}
{}
{---}






\Layout{0}{Nachgeburtsphase}
{%
Nach \VonBis{10}{20} Minuten kommt es neuerlich zu Wehen,
welche den Presswehen ähnlich sind. Es kommt zur Geburt der
\E{Plazenta} (\E{Mutterkuchen}). Die Gefahr hierbei ist,
dass es zu starkem Blutverlust bei der Mutter kommen kann.
\EE{Die Plazenta ist für weitere
Untersuchungen\footnote{Kontrolle der Vollständigkeit,
Prüfen auf Insuffizienzzeichen, \ldots} ins
Krankenhaus mitzunehmen!} 
}{
\begin{itemize}  
    \item Nach \VonBis{10}{20} Min. neuerlich Wehen. 
    \item Geburt der Plazenta.
\end{itemize}
}{}{}{}



\Layout{2}{\TxLink}{%
\begin{itemize}
  \item Birth of Baby (Vaginal Childbirth):
\url{http://www.youtube.com/watch?v=Xath6kOf0NE}
\end{itemize}

}{}{}{}{}


\begin{PictureSeriesWide}{}{Nabelschnur}
\PictureSeriesRow{3}
{./images/wikipedia-pd/Newborn_umbilical_suction-AASS-0030mm.jpg}{Durchtrennung der Nabelschnur, die patientenferne Klemme ist durch die Hand verdeckt. \SourcePicture{WMC: \enquote{Jeremykemp} / PD}}
{./images/wikipedia-pd/Navelklem.jpg}{Abgeklemmte Nabelschnur \SourcePicture{WMC: \enquote{Harmid} / PD}}
{./images/wikipedia-pd/Umbilicalstump.jpg}{Nabelschnurstumpf nach sieben Tagen \SourcePicture{WMC: \enquote{Fatrabbit} / PD}}
{}{}
\end{PictureSeriesWide}



\subsubsection{Die bevorstehende Geburt}


\Layout{0}{Anamnese -- Spezielle Fragen}
{%
\begin{itemize}
  \item \Speech{\enquote{In der wievielten SSW befinden Sie
  sich? Wann ist der geplante Geburtstermin?}}
  \item \Speech{\enquote{Die wievielte Geburt ist das?}} Bei
  Mehrgebärenden ist der Geburtsverlauf oft rascher.
  \item \Speech{\enquote{In welchem Abstand sind die Wehen?}} Je
  nachdem ist erkennbar in welcher Geburtsphase sich die
  Gebärende befindet
  \item \Speech{\enquote{Gab es schon einen Blasensprung?}} Wenn ja
  darf die Gebärende nur mehr liegend transportiert werden!
  Gefahr: Nabelschnurvorfall!
  \item \Speech{\enquote{Gab es während der Schwangerschaft
  Komplikationen?}}
  \item \T{Mutter-Kind-Pass} verlangen. Darin stehen mit
  unter wichtige Informationen über die Schwangerschaft bzw. die  Geburt.
  \item \Speech{\enquote{In welchem Krankenhaus ist die Geburt
  angemeldet?}} Sofern nicht im Mutterkindpass vermerkt.
\end{itemize}
}{
\begin{itemize}
  \item SSW, Termin?
  \item Wievielte Geburt?
  \item Wehenabstand?
  \item Blasensprung?
  \item Komplikationen bekannt?
  \item Mutter-Kind-Pass
  \item Voranmeldung zur Geburt?
\end{itemize}
}{}{}{}

\Layout{2}{Der Transport}
{%
Befindet sich die Mutter noch in der Eröffnungsphase wird
ein schnellstmöglicher, schonender Transport ins Krankenhaus
angestrebt. Grundsätzlich wird die Gebärende so
transportiert wie es für Sie am angenehmsten ist. Ab dem
erfolgtem Blasensprung darf sie jedoch nur mehr liegend
transportiert werden. Sollten Sie die Mutter stehend
angetroffen haben, wird am Transportschein dokumentiert:
\enquote{stehend angetroffen; liegend transportiert}.
Um die Gefahr eines \E{Vena-cava-Kompressionssyndroms}
(\FullPrettyRef{topic:vena-cava-syndrom}) zu vermeiden,
sollte die Patientin liegend in leichter \EE{Linksseitenlage}
transportiert werden.


}
{}
{./images/teaser/rtw-asb-921-alt.jpg}
{Ein RTW.}
{Gabriel.}



\subsubsection{Die Haus- und Wagengeburt}

\TODO{GABS}{2010-09-17}{Haus- und Wagengeburt}{}


Ab Beginn der Presswehen bzw. wenn der Kindskopf oder -teile
(Nabelschnur, Hand, \ldots) zwischen den Schamlippen
sichtbar sind, ist der weitere Transport zu
unterlassen und das Fachpersonal muss bei der Geburt
Hilfestellung leisten.

\MassnahmenMK{Spezielle
Maßnahmen}{Geburt}{m:geburt}{
}{}

\SlideCompensation{1}

\MassnahmenMK{Spezielle Maßnahmen}
{Versorgung des Neugeborenen}
{m:neugeborenes}
{
}{}

\SlideCompensation{2}

\MassnahmenMK{Spezielle Maßnahmen}
{Basisreanimation des Neugebohrenen}
{m:neugeborenes-basisreanimation}
{
}{}

\SlideCompensation{2}
% \MassnahmenNG{Spezielle Maßnahmen}{Geburt}{\label{m:geburt}}{
% Vorbereitung:
% \begin{itemize}
%   \item (Not-)Arzt hinzuziehen
%   \item Fahrzeug/Umgebung einheizen
%   \item Vorbereiten des Geburtenset
%   \item Lagerung: Z.\,B. erhöhter Oberkörper, angewinkelte
%   Beine; Wunsch der Mutter berücksichtigen!
% \end{itemize}
% Geburt:
% \begin{itemize}
%   \item Stuhlschutz
%   \item Versorgung des Neugeborenen
%   (\prettyref{m:neugeborenes})
%   \item Nachdem die Nabelschnur durchgeschnitten
%   wurde, mütterliche Seite an
%   Oberschenkel fixieren und vor Zug schützen
%   \item Dokumentation von Geburtszeit und Geburtsort
% \end{itemize}
% }
% 
% 
% \MassnahmenWide{Spezielle Maßnahmen}
% {Versorgung des Neugeborenen}
% {\label{m:neugeborenes}}
% {
% \begin{itemize}
%   \item Neugeborenes abrubbeln
%   \item Erstbeurteilung:
%   \begin{itemize}
%     \item Muskeltonus
%     \item Atemweg
%     \begin{itemize}
%       \item Bei Atemwegsverlegung: Mittels Oro-Sauger absaugen
%       \hfill\Commentary[Absaugen Neugeborener]{ASBÖ
%       Lehrmeinung 18/2011 \enquote{Absaugen Neugeborener im
%       Rettungsdienst} \cite{AsboeLehrmeinung2011-19}: \Speech{\enquote{Mehrere
%       Studien zeigten keine Verbesserung im Outcome
%       gegenüber abgesaugter Neugeborener. Daher ist es
%       obligat, keine Verzögerung bei den BLS-Maßnahmen
%       hervorzurufen, welche durch lange Absaugversuche
%       zustande kommen können. Lediglich die
%       Atemwegsobstruktion durch Mekonium muss behandelt
%       werden, indem die Atemwege freigemacht werden. Selbst
%       missfarbiges Fruchtwasser stellt keine Indikation zur
%       Absaugung
%       dar.}}
%       \ParBugfix
%       ERC 2010 \cite{Erc2010Sek06De}: \Speech{\enquote{Absaugen ist
%       nur notwendig, wenn die Atemwege verlegt sind. Eine
%       solche Verle- gung kann aufgrund von Mekonium, auch
%       wenn das Neugeborene keine Mekoniumablagerungen auf
%       der Haut zeigt, aber auch Blutkoageln, zähem Schleim
%       oder Vernix bestehen. Wird abgesaugt, ist zu bedenken,
%       dass zu heftiges oropharyngeales Ab- saugen das
%       Einsetzen einer suffizienten Spontanatmung verzögern
%       und zu einem Laryngospasmus sowie zu einer vagalen
%       Bradykardie führen kann. Das Vorhandensein von
%       zähflüssigem Mekonium beim schlaffen, avitalen
%       Neugeborenen ist die einzige Situation, in der ein
%       sofortiges Absaugen des Oropharynx zu erwägen ist.}}
%       }%Commentary
%     \end{itemize}
%     \item Atmung
%     
%     Bei Schnappatmung oder fehlender Atmung: Atemwege öffnen
%     5 initiale Beatmungen
%     \item Kreislauf: \Abk{HF} messen (durch Auskultation
%     an der Herzspitze (nur geübtes Personal) oder an
%     \E{Oberarmarterie}). Wenn ein \E{geeignetes
%     Pulsoxymeter} vorhanden ist, kann auch dieses verwendet werden.
%     \hfill\Commentary[Ort der HF-Messung beim Neugeborenen]
%     {ERC 2010 \cite{Erc2010Sek06De}: \Speech{\enquote{Die beste Methode
%     zur Beurteilung der Herzfrequenz ist die direkte Auskultati- on mit dem Stethoskop über der Herzspitze. Das Tasten des
%     Pulses an der Basis der Nabelschnur ist oft möglich, kann
%     aber irreführend sein. Eine Beurteilung der Herzfrequenz
%     allein über die Pulsation der Nabelschnur ist nur
%     zuverlässig, wenn die Herzfrequenz über 100\PerMin* liegt.}}
%     \ParBugfix
%     Bei Sanitätern, insbesondere der
%     niedrigeren Ausbildungsstufen, kann nicht davon
%     ausgegangen werden, dass eine Auskultation in dieser
%     stressreichen Situation fehlerfrei und zuverlässig angewendet werden kann. Die
%     Oberarmarterie erscheint in diesem Fall als
%     pragmatischer Kompromiss.}
%   \end{itemize}
% \end{itemize}
% {\FormatTable
% \begin{tabularx}{\linewidth}{XX}
% \toprule\addlinespace
% \multicolumn{1}{c}{\includegraphics[width=1cm]{./icons/sm-basic.png}}
% &
% \multicolumn{1}{c}{\includegraphics[width=1cm]{./icons/sm-pale.png}~~\includegraphics[width=1cm]{./icons/sm-pale-frown.png}}
% \\\addlinespace
% \noindent\TabHeading{Kräftiges Schreien/suffiziente Atmung}
% &
% \noindent\TabHeading{Insuffiziente Spontanatmung,
% Atemstillstand oder Bradykardie}
% \\\addlinespace\cmidrule(lr){1-1}\cmidrule(lr){2-2}\addlinespace
% \begin{itemize}
%   \item[\eye] Guter Muskeltonus,
%   \item[\eye] Herzfrequenz \textgreater\,100\PerMin*
% \end{itemize}
% 
% 
% 
% 
% &
% \begin{itemize}
%   \item[\eye] Normaler, reduzierter Muskeltonus, oder sogar
%   schlaffer Muskeltonus (\enquote{floppy})
%   \item[\eye] Herzfrequenz \textless\,100\PerMin*; Bradykarde
%   oder nichtnachweisbare Herzfrequenz,
%   \item[\eye] Oft mit ausgeprägter Blässe als Zeichen
%   einer schlechten Perfusion.
% \end{itemize}
% 
% \\\addlinespace\cmidrule(lr){1-1}\cmidrule(lr){2-2}\addlinespace
% \begin{itemize}
%   \item Abtrocknen und Einwickeln in warme Tücher:Warm
%   einpacken (nur das Gesicht darf frei bleiben) und der Mutter
%   überreichen.
%   \item Abnabeln, frühestens nach 1\Min*:
%   \begin{enumerate}
%     \item Erst wird eine Handbreite vom Kind
%     entfernt die erste Nabelklemme gesetzt.
%     \item Dann wird die Nabelschnur Richtung Mutter ausgestreift und
%     wieder eine Handbreit entfernt die zweite Klemme gesetzt.
%     \item Zwischen den beiden
%     Klemmen wird die Nabelschnur durchtrennt.
%   \end{enumerate}
%   \item Das Neugeborene
%   kann der Mutter auf den Bauch gelegt werden. Durch den
%   Kontakt zur Haut der zugedeckten Mutter wird das Baby
%   gewärmt.
%   
%   \item Es kann zu diesem Zeitpunkt bereits an die Brust
%   angelegt werden.
% \end{itemize}
%   &
%   % Nachdem ein Neugeborenes dieser Gruppe getrocknet und in
%   % warme Tücher gewickelt wurde, ist meist eine kurze
%   % Maskenbeatmung ausreichend. Kommt es darunter nicht zu einem
%   % adäquaten Anstieg der Herzfrequenz, müssen Herzdruckmassagen
%   % durchgeführt werden.
%   
%   \begin{itemize}
%     %     \item \ce{O2}-Gabe (\enquote{\ce{O2}-Dusche}) für 30\,Sek.
%     \item Atemwege freimachen, absaugen
%     \item 5 Beatmungen
%     \item Wenn keine Steigerung der Herzfrequenz:
%     \begin{itemize}
%       %     \item \ce{O2}-Gabe (\enquote{\ce{O2}-Dusche}) für 30\,Sek.
%       \item Lagerung und Atemwege überprüfen
%       \item 5 Beatmungen wiederholen
%     \end{itemize}
%     \item Wenn weiterhin keine ausreichende Atmung oder
%     \EE{Herzfrequenz \textless\,60\PerMin*}:
%     \begin{itemize}
%       \item \Sign{→} Reanimation des Neugeborenen
%     \end{itemize}
%     \FullPrettyRef{m:neugeborenes-basisreanimation}
%   \end{itemize}
%   
%   % Nachdem ein Neugeborenes dieser Gruppe getrocknet und in
%   % warme Tücher gewickelt wurde, müssen unverzüglich die
%   % Atemwege geöffnet werden. Die Lungen müssen belüftet werden,
%   % und das Kind muss beatmet werden. Wurden die Atemwege
%   % erfolgreich geöffnet und die Lungen belüftet, können
%   % außerdem Herzdruckmassagen und möglicherweise eine
%   % Medikamentengabe notwendig sein. Eine sehr kleine Gruppe von
%   % Neugeborenen bleibt trotz adäquater Spontanatmung und guter
%   % Herzfrequenz hypox­ ämisch. Die Diagnosen umfassen hier ein
%   % weites Spektrum, wie z. B. eine Zwerchfellhernie,
%   % Surfactant-Mangel, eine kongenitale Pneumonie, einen
%   % Pneumothorax oder ein angeborenes zyanotisches Herzvitium.
%   
%   
%   \\\addlinespace\bottomrule
% \end{tabularx}
% }
% 
% \cite{Erc2010Sek06De,AsboeLehrmeinungRd2011-19,AsboeLehrmeinungRd2011-18}
% }
% 
% \MassnahmenNG{Spezielle Maßnahmen}
% {Basisreanimation des Neugebohrenen}
% {\label{m:neugeborenes-basisreanimation}}
% {
% \begin{itemize}
%   \item 2 Daumen nebeneinander über dem
%   unteren Drittel des Brustbeins, direkt unter einer gedachten Linie zwischen den
%   Brustwarzen
%   
%   \item Umgreifen des gesamten
%   Brustkorbs und Unterstützung des Rückens des Kindes.
%   \item Thoraxkompressionen und Beatmungen im Verhältnis von
%   \EE{3\,:\,1}
%   
% %   \item Ziel: 90 Kompressionen und 30 Beatmungen\PerMin*
% %   (120 Maßnahmen\PerMin*)
%   \item \TxMassVit
% \end{itemize}
% 
% \cite{Erc2010Sek06De}
% }

Nach dem Ende der Geburt müssen die \EE{Geburtszeit} und der
\EE{Geburtsort} (Straße mit Hausnummer; Kreuzung; Autobahn
mit km-Angabe, etc.) \EE{dokumentiert} werden. Der
Geburtsort ist insofern wichtig, da für die Ausstellung der
Geburtsurkunde  das für den Geburtsort zuständige Bezirksamt
verantwortlich ist. Wenn Mutter und Kind wohlauf sind, kann
der Transport in ein Krankenhaus fortgesetzt werden.

Es soll nicht auf die Nachgeburt gewartet werden. Sollte die
Plazenta vor Eintreffen im Krankenhaus geboren werden, ist
diese unbedingt sicherzustellen und mitzunehmen. Sie wird im
Krankenhaus auf Vollständigkeit und Reife untersucht.



\subsection{Notfälle während der Geburt}
\label{topic:notfall-geburt}

\subsubsection{Nabelschnurvorfall}

Bei einem Nabelschnurvorfall in den Geburtskanal bestehen
zwei Gefahren:
Einerseits wird während des Geburtsvorganges die Nabelschnur
durch den Kindskopf abgeklemmt, und somit ist die
Sauerstoffversorgung des Kindes nicht mehr gewährleistet.
Andererseits kann sich die Nabelschnur  um den Hals des
Kindes wickeln und es kommt während der Geburt zur
Strangulation des Kindes.

\MassnahmenMK{Spezielle Maßnahmen}{Nabelschnurvorfall}
{m:nabelschnurvorfall}
{

}
% \MassnahmenNG{Spezielle Maßnahmen}{Nabelschnurvorfall}
% {\label{m:nabelschnurvorfall}}
% {
% \begin{itemize}
% 	\item Beckenhochlagerung
% 	\item Mutter soll \underline{nicht} pressen (wenn möglich Geburt erst im
% 	Krankenhaus)
% 	\item Wehenhemmung (hecheln; \Z{medikamentös durch Notarzt})
% 	\item Nabelschnurumschlingung des kindlichen Halses während der Geburt sofort lösen
% \end{itemize}
% }

\subsubsection{Placenta praevia}

\Layout{0}{\TxBeschreibung}{%
Normalerweise liegt die Plazenta nicht vor dem Geburtskanal.
Bei einer sehr tiefen Einnistung des befruchteten Eis in der
Gebärmutter, kann es jedoch dazukommen, dass sich die
\EE{Plazenta \E{vor} dem Geburtskanal} bildet. Dies ist
natürlich ein sehr großes Geburtshindernis und i.\,d.\,R.
ist eine vaginale Geburt unmöglich. Außerdem kann es durch
vorzeitiges Lösen des Mutterkuchens zu massiven Blutungen
kommen.
Eine Placenta praevia ist meistens schon durch die
\E{Kontrolluntersuchungen} beim Frauenarzt bekannt und kann
erfragt bzw. im Mutter-Kind-Pass nachgelesen werden.
}{%
\begin{itemize}
  \item Plazenta liegt vor dem Geburtskanal
  \item Vaginale Geburt nicht möglich
  \item Blutungsgefahr 
  \item Anamnese, Mutter-Kind-Pass
\end{itemize}
}{}{}{}


\Layout{0}{\TxSymptome}{%
Unter Wehentätigkeit kommt es zu starken Blutungen aus der
Scheide.

\vspace{0.5cm}

}{%
\begin{itemize}
  \item Vaginalblutungen während Wehentätigkeit
\end{itemize}
}{}{}{}

\MassnahmenMK{Spezielle Maßnahmen}{Placenta praevia und
bevorstehende Geburt}{m:placenta-praevia}
{
}{}
% \MassnahmenNG{Spezielle Maßnahmen}{Placenta praevia und bevorstehende
% Geburt}{\label{m:placenta-praevia}}
% {
% \begin{itemize}
%     \item \TxMassVitBes
%     \begin{itemize}
%         \item Lagerung: Sterile Vorlage und Lagerung nach
%         Fritsch.
%         \item Ggfs. Schockbekämpfung
%     \end{itemize}
% \end{itemize}
% }

% \paragraph{Gefahren:}
% 
% Abklemmung der Nabelschnur  (Sauerstoffversorgung  [F0EA?] )
% 
% Strangulation (Gehirnperfusion [F0EA?] )
% 
% \paragraph{Therapie:}
% 
% Linksseitenlagerung, Beckenhochlagerung
% 
% Mutter soll NICHT pressen (Geburt w.m. erst  im KH!!!)
% 
% Wehenhemmung (hecheln, medikamentös) 
% 
% Nabelschnurumschlingung des  kindlichen Halses während der Geburt 
% sofort lösen!
% 
% Placenta praevia
% 
% Plazenta liegt unterhalb des Fetus in  Geburtskanal
% 
% Geburtshindernis
% 
% massive Blutungen
% 
% \paragraph{Maßnahmen}
% 
% Lagerung nach Fritsch
% 
% Schockbekämpfung
% 
% NAW
% 
% im KH: Notsectio

\subsubsection{Pathologische Geburtslagen}


Alle Lagen des Kindes, bei denen nicht zuerst der Kopf im Geburtskanal erscheint.
\begin{itemize}
	\item Steißlage bzw. Beckenendlage
	\item Armvorfall
\end{itemize}

Die größte Gefahr bei allen Lageanomalien ist, dass die
Nabelschnur abgedrückt wird. Daher sollte unter allen
Umständen die Geburt im KH, mittels Kaiserschnitt,
durchgeführt werden. Die Geburt in jedem Fall verhindern
bzw. verzögern. Ein Armvorfall ist eine geburtsunmögliche
Lageanomalie. 

\MassnahmenMK
{Spezielle Maßnahmen}
{Pathologische
Geburtslagen}
{m:pathologische-geburtslagen}
{}{}


\subsubsection{Uterusatonie}

\Layout{0}
{\TxBeschreibung}
{Normalerweise zieht sich in der Nachgeburtsphase die
Uterusmuskulatur zusammen. In diesem Fall kommt es aber zur
Erschlaffung (Atonie) der Gebärmutter. Die Folge sind
lebensbedrohliche Blutungen durch das fehlende Kontrahieren
des Uterus nach der Plazentalösung.}
{
\begin{itemize}
  \item Erschlaffung der Gebärmutter
  \item Blutungen
\end{itemize}
}
{}
{}
{}


\MassnahmenMK{Spezielle
Maßnahmen}{Uterusatonie}{m:uterusatonie}{
}{}
% \MassnahmenNG{Spezielle
% Maßnahmen}{Uterusatonie}{\label{m:uterusatonie}}{
% \begin{itemize}
%   \item \TxMassVitBes
%   \begin{itemize}
%     \item Schockbekämpfung
%     \item Lagerung: Fritsch, Beine hoch
%   \end{itemize}
% \end{itemize}
% }

% normalerweise zieht sich in der  Nachgeburtsphase die Uterusmuskulatur 
% zusammen
% 
% hier: Erschlaffung der Uterusmuskulatur  nach Plazentalösung
% ={\textgreater} großflächige  Blutung!
% 
% \paragraph{Maßnahmen}
% 
% \begin{itemize}
%     \item Allgemeine Maßnahmen bei lebensbedrohlichen Zuständen
%     \item Schockbekämpfung
%     \item Lagerung: Fritsch, Beine hoch
%     \item Evtl. Druck mit beiden Fäusten oberhalb  Schambein ${\rightarrow}$
%     Blutungsverminderung
%     \item Zügiger Transport mit Voranmeldung
% \end{itemize}


\subsubsection{Uterusruptur}

\Layout{2}{\TxBeschreibung}{% 
Die Uterusruptur kommt meistens nur nach vorangegangenem
Kaiserschnitt vor. Es treten während der Wehen plötzlich
Schmerzen auf, welche anders als Wehenschmerzen empfunden
werden. Die Wehen nehmen dann ab oder hören plötzlich auf.
Damit wird der Geburtsvorgang gestoppt und es herrscht
Lebensgefahr für die Mutter aufgrund starker Blutungen und
für das Kind wegen des \ce{O2}-Mangels.
}{}{}{}{}

\MassnahmenMK{Spezielle
Maßnahmen}{Uterusruptur}{m:uterusruptur}{
}{}
% \MassnahmenNG{Spezielle
% Maßnahmen}{Uterusruptur}{\label{m:uterusruptur}}{
% \begin{itemize}
% 	\item \TxMassVit
% 	\item Schockbekämpfung
% 	\item Lagerung: Fritsch, Beine hoch
% \end{itemize}
% }

% 
% {\mdseries
% plötzlicher Schmerz während Wehen,  dann keine Wehen mehr
% ${\rightarrow}$ Geburtsstop}
% 
% \Fazit{Lebensgefahr für Mutter (Blutung) und  Kind (O2
% -Mangel!)}
% 
% \paragraph{Therapie}
% 
% \begin{itemize}
%     \item Allgemeine Maßnahmen bei lebensbedrohlichen Zuständen
%     \item Schockbekämpfung
%     \item Zügiger Transport mit Voranmeldung
% \end{itemize}


\subsubsection{Asphyxie während der Geburt}

\Definition{Sauerstoffunterversorgung (Hypoxie) des
Neugeborenen infolge einer atem- oder kreislaufbedingten
Störung vor, während oder nach der Geburt.}

\Layout{2}{\TxBeschreibung}{%
Die möglichen Ursachen für die Hypoxie sind vielfältig und
umfassen u.\,a. eine Placenta praevia, einen
Nabelschnurvorfall, Lage\-ano\-malien, eine Frühgeburt,
u.\,v.\,a.\,m. Wird die Ursache nicht rechtzeitig behoben
und bleibt die Hypoxie bestehen, kommt es zu einer
irreperablen Hirnschädigung.
}{%

}{}{}{}

\Layout{0}{\TxSymptome}{%
Das Leitsymptom einer Asphyxie ist eine \EE{Herzfrequenz <
100\PerMin*}, eine \EE{unzureichende Atmung}
(Schnappatmung, fehlende Atmung) und/oder ein \EE{schlaffer
Muskeltonus} des Kindes
Weitere Hinweise auf  eine Asphyxie sind eine grünlich-trübe
Verfärbung des Fruchtwassers sowie eine \E{fortbestehende
Zyanose} (eine Zyanose unmittelbar nach der Geburt ist
normal).
}{%
\begin{itemize}
  \item HF < 100\PerMin*
  \item Unzureichende Atmung
  \item Schlaffer Muskeltonus
\end{itemize}
}{}{}{}

\MassnahmenMK{Spezielle Maßnahmen}{Asphyxie des
Neugeborenen}
{m:asphyxie}
{%
Die Behandlung der Asphyxie erfolgt im Rahmen der Versorgung
des Neugeborenen (\FullPrettyRef{MK-m:neugeborenes}).
}{}
% \MassnahmenNG{Spezielle Maßnahmen}{Asphyxie des
% Neugeborenen}
% {\label{m:asphyxie}}
% {%
% Die Behandlung der Asphyxie erfolgt im Rahmen der Versorgung
% des Neugeborenen (\FullPrettyRef{m:neugeborenes}).
% % \begin{itemize}
% %   \item \TxMassVit
% %   \item Warmhalten
% %   \item Abtrocknen/abrubbeln; stimulieren: Fußsohlen und
% %   Rücken klopfen, reiben
% %   \item Atmung und Herzfrequenz beurteilen
% %   \item Bei Schnappatmung, fehlender Atmung oder HF <
% %   100\PerMin*:  \begin{itemize}
% % %     \item \ce{O2}-Gabe (\enquote{\ce{O2}-Dusche}) für 30\,Sek.
% %     \item Atemwege freimachen, absaugen
% %     \item 5 Beatmungen
% %   \end{itemize}
% %   \item Wenn keine Steigerung der Herzfrequenz: 
% %   \begin{itemize}
% % %     \item \ce{O2}-Gabe (\enquote{\ce{O2}-Dusche}) für 30\,Sek.
% %     \item Lagerung und Atemwege überprüfen 
% %     \item 5 Beatmungen wiederholen
% %     \item \ce{O2}-Gabe
% %     \end{itemize}
% %   \item Wenn weiterhin keine ausreichende Atmung oder
% %   \EE{Herzfrequenz < 60\PerMin*} 
% %   
% %   \Sign{→} Reanimation des Neugeborenen
% % \end{itemize}
% }

% \clearpage % TEMP temp clearpage
\section{Sonstige gynäkologische Erkrankungen und
Notfälle}

\subsection{Vaginale Blutungen}
\label{topic:unterleibsblutungen}

\Layout{0}{\TxBeschreibung}{%
Die vaginale Blutung ist präklinisch schwer abzuklären und
kann viele Ursachen haben. Wichtig ist festzuhalten, dass
jede vaginale Blutung abgeklärt gehört, da sich dahinter
auch ein Tumorgeschehen (besonders bei Patientinnen in der
Menopause) verstecken kann. Darüber hinaus kommt eine
verstärkte Monatsblutung, Verletzung (nach Sturz,
Geschlechtsverkehr oder nach Vergewaltigungen),
Defloration\footnote{Entjungferung} oder Blutungen im Rahmen
einer Schwangerschaft in Betracht. Bei vaginalen Blutungen
in Kindern darf man \E{nicht} automatisch von einem
erfolgten Missbrauch ausgehen.
}{
\begin{itemize}
    \item Verstärkte Monatsblutung
    \item Verletzungen (Sturz, Geschlechtsverkehr, Vergewaltigung,
    \ldots)
    \item Defloration
    \item Tumor
    \item Im Rahmen einer Schwangerschaft
\end{itemize}
}{}{}{}

Eine akute vitale Bedrohung besteht seltenst, dennoch muss
sie routinemäßig vor allem anhand des Blutverlustes und der
Vitalwerte eingeschätzt werden. 

\MassnahmenMK{Spezielle Maßnahmen}{Vaginaler
Blutung}{m:vaginale-blutung}{
}{}
% \MassnahmenNG{Spezielle Maßnahmen}{Vaginaler
% Blutung}{\label{m:vaginale-blutung}}{
% \begin{itemize}
%     \item Lagerung nach Fritsch
%     \item Transportentscheidung: Abteilung für
%     Gynäkologie, bei Wehen Kreißsaal
% \end{itemize}
% }


\begin{Literary}
  {\fullcite{HaagGyn2007}} 

  {\fullcite{BuehlingGynaekologie1}}

  {\fullcite{OHCSP7}}

  {\fullcite{Nilsson1997}} 
  
  Weiters: \cite{
SiewertChirurgie8,
AMLS3,
% Erc2005DeSec5,
% Erc2005Sec5,
% Erc2005DeSec6,
% Erc2005DeSec3,
% Erc2005Sec3,
BoehmerTaschRett6,
Fauci2008,
LPN3,
Boecker2,
% Helfen2008,
LippertAnatomie7,
RedelsteinerHRNS1,
ForMed2,
StriebelAn7,
KlinkePhysio3,
DRAn3,
ForthPharma8,
LuellmannPharma16,
AnCompact3,
RueckerBildatlas1}
\end{Literary}

% \section*{\protect\dsliterary{} Weiterführende Literatur}
% \begin{WidePar}
% \begin{multicols}{2}
% \Z{%
% \begin{labeling}{mmm}
%   \item[\cite{HaagGyn2007}]
%   \emph{\bibentry{HaagGyn2007}} 
%   \item[\cite{BuehlingGynaekologie1}] 
%   \emph{\bibentry{BuehlingGynaekologie1}}
%     \item[\cite{OHCSP7}] 
%   \emph{\bibentry{OHCSP7}}
%   \item[\cite{Nilsson1997}] 
%   \emph{\bibentry{Nilsson1997}} 
% \end{labeling}
% }
% \end{multicols}
% \end{WidePar}



\nocite{Juurlink2005,
Kalla2006,
Vlcek2008,
Wells2000,
Patzelt2005,
Sachsenweger2003,
}

\ChapterBibliography
